\documentclass[11pt]{article}
\usepackage[margin=2.5cm]{geometry}
\usepackage{amsmath}
\usepackage{changepage}
\usepackage{xcolor}
\usepackage[most]{tcolorbox}
% the framed fixture's solved box, drawn by tcolorbox the way mdframed draws it there: a 4pt rule on the left, 35pt
% out in the margin, the text 30pt to its right, no other frame
\tcbset{frametitle/.style={title={#1}}}
\newtcolorbox{solved}[1][]{
  breakable, enhanced, parbox=false, frame hidden, boxrule=0pt, boxsep=0pt, arc=0pt,
  colback=white, borderline west={4pt}{0pt}{teal},
  enlarge left by=-35pt, width=\linewidth+35pt,
  left=34pt, right=0pt, top=6pt, bottom=6pt,
  before skip=12pt, after skip=12pt,
  fonttitle=\bfseries, coltitle=black, colbacktitle=white, titlerule=0pt,
  toptitle=6pt, bottomtitle=0pt, lefttitle=34pt,
  #1
}
\begin{document}

\begin{adjustwidth*}{0mm}{}

\section*{Wrapped document}

Margin paragraph marginone sits inside a margin wrapper that spans the whole document, the
way some course notes put every page in one adjustwidth environment, and it wraps the
column a few times so the layout has something to measure.

Margin paragraph margintwo follows it with a different word count so the two paragraphs
break their lines in different places on the page.

\begin{solved}[frametitle={Solved example}]
Framed paragraph frameone opens the solved example inside a coloured frame whose rule runs
down the left side, and it goes on long enough to wrap the frame width two or three times.

Framed paragraph frametwo continues the example with an inline formula $V_o = -V_i R_2/R_1$
and a sentence that keeps going so that it, too, wraps inside the frame.
\begin{equation*}
  A = -\frac{R_2}{R_1}
\end{equation*}
Framed paragraph framethree comes after the display, still inside the frame, and closes the
example with a last sentence that wraps once.
\end{solved}

Margin paragraph marginthree comes after the frame and before the end of the wrapper, a plain
paragraph again at the full text width.

\end{adjustwidth*}

\end{document}
